\section{The CI Labels and Policy Gate}
\label{sec:system_arch}

\subsection{The CI Label}
\label{sec:label}
A ruling should be explainable: although the gate decides each flow at request
time, its decision is predictable from the label and policy written in advance.
The label therefore encodes the parameters of an information
flow~\cite{nissenbaum}. Recipient comes from the requester's credentials; Data
Subject, Information Type, and Transmission Principle come from the document
label; Sender is implicit, as it is always the Store.

\begin{table}[t]
\centering
\caption{Fields of the document label that constrain each Contextual Integrity
parameter.}
\label{tab:ci-mapping}
\footnotesize
\resizebox{0.82\columnwidth}{!}{%
\begin{tabularx}{\columnwidth}{@{}lX@{}}
\toprule
\textbf{CI parameter} & \textbf{Constraining label field(s)} \\
\midrule
Data subject           & \texttt{subject\_type} \\
Sender                 & \emph{Implicit (always the Store)} \\
Recipient              & \texttt{permitted\_receiver\_roles} \\
Information type       & \texttt{sensitivity\_tier}, \texttt{pii\_attributes},
                         \texttt{implied\_attributes}, \texttt{extra\_sensitive},
                         \texttt{document\_type} \\
Transmission principle & \texttt{governing\_principle} \\
\bottomrule
\end{tabularx}}
\end{table}

Each parameter serves auditability, decision making, or both. Data Subject
defines who the document is about, and Recipient establishes the other half of
the relationship. Information Type determines how tightly that relationship
must be established: for sensitive tiers a permitted role is not sufficient,
and the requester must also hold an authorized relationship with the subject.
The transmission principle is not evaluated as a separate condition; the role
and relationship requirements implement it, and \texttt{governing\_principle}
records which principle they encode.

While Contextual Integrity is granular about individuals, the label is written
over \textit{classes}: the Recipient field could hold ``Doctor'' but never a
specific physician. We call this field, \texttt{permitted\_receiver\_roles},
the document's \textit{access-control list} (ACL). Individual-level
authorization is carried by the assignment relation (Section~\ref{sec:gate}).
Labels attach to whole documents, so a record containing both a diagnosis and
a billing code is treated indivisibly.

Not every field bears on the ruling. Table~\ref{tab:label-fields} separates the
fields the gate decides from, those it reports in its justification, and those
retained for provenance. The two deciding fields differ in kind: the tier is
descriptive, and the document is evidence for it; the ACL is normative.

\begin{table}[t]
\centering
\caption{Fields of the document label by their role in the gate's decision
procedure.}
\label{tab:label-fields}
\footnotesize
\resizebox{0.82\columnwidth}{!}{%
\begin{tabularx}{\columnwidth}{@{}l>{\raggedright\arraybackslash}X@{}}
\toprule
\textbf{Role in the gate} & \textbf{Label fields} \\
\midrule
Decides  & \texttt{permitted\_receiver\_roles}, \texttt{sensitivity\_tier} \\
\addlinespace
Explains & \texttt{document\_type}, \texttt{extra\_sensitive}, \texttt{governing\_principle} \\
\addlinespace
Records  & \texttt{originating\_context}, \texttt{subject\_type}, \texttt{pii\_attributes},
           \texttt{implied\_attributes}, \texttt{extra\_sensitive\_reason},
           \texttt{assignment\_required} \\
\bottomrule
\end{tabularx}
}
\end{table}

\subsection{The CI Policy Gate}
\label{sec:gate}
The gate receives a request from the Assistant, resolves it to a set of
candidate documents (Section~\ref{sec:retrievalpath}), and evaluates each as a
potential flow. Each evaluation is a function of six inputs: the requester's
role, that role's context, the requester's assignment set, and the document's
ACL, tier, and subject. It returns a decision for each candidate and the rule
that produced it.

A role is authoritative only in the context that grants it. This follows from
the (Agent, Role) pairing in the CI formalism~\cite{barth2006}: the same
individual acting in a different capacity is a different principal. A Doctor
acting in the administrative context is refused, because the Doctor role is
authoritative only in the medical context.

Beyond document labels, the {\em Policy Authority} maintains an
\textit{assignment relation}: the subjects each requester holds an authorized
relationship with. Requiring both a permitted role and an assignment separates
this from ordinary role-based access control~\cite{sandhu1996}, under which any
Doctor could access any patient's chart. The role establishes the scope of
access; the assignment establishes need-to-know. A requester with no
assignments can reach only documents whose tier does not require one
(Table~\ref{tab:tiers}).

\begin{table}[t]
\centering
\caption{Sensitivity tiers and whether each requires an assignment.
\texttt{public} and \texttt{internal} documents concern the organization, so
there is no subject to be assigned.}
\label{tab:tiers}
\footnotesize
\resizebox{0.80\columnwidth}{!}{%
\begin{tabularx}{\columnwidth}{@{}l>{\raggedright\arraybackslash}Xc@{}}
\toprule
\textbf{Tier} & \textbf{Content} & \textbf{Assignment} \\
\midrule
\texttt{public}              & Freely shareable; no private content                      & No  \\
\texttt{internal}            & Organizational; not about an individual                   & No  \\
\addlinespace
\texttt{pii\_basic}          & Identifiers only, no clinical or financial content        & Yes \\
\texttt{clinical\_standard}  & Diagnoses, medications, clinical history                  & Yes \\
\texttt{clinical\_sensitive} & Mental health, substance use, HIV/AIDS, reproductive care & Yes \\
\texttt{financial}           & Billing codes, insurance numbers, salary figures          & Yes \\
\texttt{employment}          & Contracts, HR records, reviews, credentialing             & Yes \\
\bottomrule
\end{tabularx}
}
\end{table}

The gate returns \textsc{Permit} only when every check holds; \textsc{Deny} is
the default. On a denial, the system can either withhold the document silently
or tell the model why.

The decision is a function of these six inputs and nothing else: it does not
consult the request text, the Assistant's reasoning, or any prior request, so
identical requests yield identical decisions indefinitely (verified in
Section~\ref{sec:conformance}). Appendix~\ref{app:alg} gives the decision
procedure.